\origin{ai}
\subsection{PRDM9 Hotspot Rows as a Boundary Test for a Translation-Conditioned RNA Half-Life Bridge}
\label{sec:bioreality-prdm9-hotspot-q6-rbp-bridge}

\paragraph{Finite empirical row set.}
The finite contact in this record is the PRDM9 meiotic recombination hotspot panel derived from the Pratto 2014 surface.  The parsed run contains $52963$ hotspot rows and completed in $168.443$ seconds under run id $e57a09da690c$.  Those rows give a concrete DNA-side cross-layer measurement: each row belongs to a hotspot catalogue whose biological interpretation is tied to PRDM9 motif preference and double-strand-break placement.  The object tested here is not whether PRDM9 participates in meiotic hotspot localization; the row count records that a finite hotspot surface is present.  The question is whether such a surface can be used as support for the separate conjecture that $B^{*}_{Q6}$ synonymous residual coordinate windows in coding transcripts carry a translation-conditioned RNA half-life effect through a measured decay-factor RBP bridge.  It cannot do so by itself, because the PRDM9 rows live on a DNA recombination surface and do not contain matched transcript/CDS rows, same-row translation readouts, RNA half-life estimates, local decay-factor RBP occupancy, target-RBP loss, wild-type rescue, or binding-defective control rescue.

\paragraph{Required rows for the RNA half-life bridge.}
A positive RNA half-life RBP bridge would require a joined finite table whose row $i$ contains a synonymous residual coordinate label $q_i$, a same-row translation readout $T_i$, a transcript abundance or transcription control vector $A_i$, an RNA half-life readout $H_i$, a measured local occupancy value $O_i$ for a prespecified decay-factor RBP, and perturbation/rescue indicators $L_i,W_i,D_i,N_i$ for target-RBP loss, wild-type rescue, binding-defective rescue, and non-target rescue.  The admissible claim is bounded by the rowwise statistic
$$
\begin{aligned}
\Delta_H(q) &= H_i - \widehat H_i(A_i,T_i,\text{structure}_i,\text{composition}_i,\text{isoform}_i,\text{batch}_i),\\
\Delta_O(q) &= O_i - \widehat O_i(\text{motif density}_i,\text{structure}_i,\text{composition}_i,\text{batch}_i),\\
\mathcal B(q) &= \big(\Delta_H(q),\Delta_O(q),L_i,W_i,D_i,N_i\big).
\end{aligned}
$$
The bridge is present only if the half-life residual follows an experimentally separated occupancy change, collapses when the target RBP is lost, returns under wild-type rescue, and fails under binding-defective or non-target rescue.  The PRDM9 hotspot table supplies none of these RNA-window fields.  It therefore remains an informative recombination contact, not evidence for the transcript-decay mechanism demanded by this conjecture.

\paragraph{Coordinate distinction and null model.}
The relevant coordinate contrast is high-coordinate versus low-coordinate synonymous residual label after quotienting by amino-acid sequence.  Under the local null, after conditioning on transcription rate, RNA structure, motif density, GC and dinucleotide composition, codon-pair context, splice or isoform state, transcript abundance, amino-acid composition, batch, and the same-row translation readout, the high- and low-coordinate labels are exchangeable with respect to RNA half-life.  Written at the finite-row level, the null is
$$
\begin{aligned}
H_i \perp q_i \mid (T_i,A_i,\text{structure}_i,\text{composition}_i,\text{isoform}_i,\text{batch}_i),
\end{aligned}
$$
with the additional bridge null that a half-life contrast is not an RBP-mediated contrast unless the occupancy/loss/rescue pattern is present in the same finite panel.  A DNA hotspot row may be biologically real and still be irrelevant to this null, because it does not measure transcript half-life or decay-factor binding at an editable RNA window.

\paragraph{Auxiliary finite contacts and why they still do not close the bridge.}
Two other verified finite surfaces clarify the boundary.  The RBP binding contact has $900$ peak rows in its parsed panel, with $3$ cross-RBP concordant outcomes and runtime $223.648043$ seconds under run id $b225890251c2$.  This shows that finite RBP-binding rows can exist in the project, but those rows are not the required same-organism, same-transcript, same-window decay-factor occupancy/loss/rescue panel for the $B^{*}_{Q6}$ RNA half-life conjecture.  Similarly, the mTOR translation contact contains $6494$ readout rows from the GSE36892 source and ran in $9.469$ seconds under run id $0c2bbfd3e2d9$.  That translation surface can bound translation readout in its own setting, but it is not an RNA half-life measurement and does not provide target-RBP perturbation or rescue.  These finite contacts demonstrate that the required row types are conceptually separable: recombination hotspots, RBP-binding peaks, and translation readouts are different measurement layers unless a single matched panel explicitly joins them.

\paragraph{Readback verdict.}
For this standalone chapter, the finite conclusion is a guarded negative bridge verdict.  The PRDM9 hotspot panel is a valid finite recombination surface with $52963$ hotspot rows, but it cannot be promoted into a $B^{*}_{Q6}$ translation-conditioned RNA half-life RBP bridge.  The conjecture would become testable only after adding matched transcript/CDS rows carrying synonymous residual coordinate labels, same-row translation readouts, transcript abundance controls, RNA half-life measurements, measured decay-factor RBP occupancy at the relevant RNA window, target-RBP loss, wild-type rescue, and binding-defective or non-target rescue controls.  Without that joined table, the strongest local statement is that the PRDM9 contact records DNA-side hotspot structure, while the RNA half-life bridge remains unclosed.

\paragraph{Local cannot-claim boundary.}
No translation realization follows from RNA half-life alone, and no RNA-decay mechanism follows from a translation contact name.  No RBP mechanism follows from a half-life association unless measured local occupancy, loss collapse, wild-type rescue, and failed binding-defective or non-target rescue are present in the same matched panel.  No protein output, folding, physical admissibility, molecular function, phenotype, adaptation, causation, or universal biological law follows from the PRDM9 hotspot rows, from $B^{*}_{Q6}$ coordinate geometry, from RBP peak counts, or from translation-readout routing unless a separate layer-matched measurement supplies that specific readout.